\subsection{QC-009: exact determinant and padded samples with second tensor equal to the identity}

\paragraph{Statement QC-009.}
The standard-library generator
\texttt{generate\_orthonormal\_samples.py} supplies two reproducible
sets of centered nonic tensor values modulo $101$, all with normalized
second tensor $I_3$ and with every tensor entry of orders $2$ through
$9$ retained:
\[
\begin{array}{c|r|r|r}
\text{source locus}&\text{points}&\text{coordinates per point}&
 \text{consistency checks}\\\hline
D_{9,4}&2420&216&7260\\
V=\{\ell q_8\}&2410&216&7230
\end{array}
\tag{QC-009.1}
\]
The determinant values come from explicit integer matrix pencils.
The padded values have explicit integer polynomial lifts with an
exact linear factor. The identity condition on the second tensor
is asserted modulo $101$; no claim that an integer congruence
normalizes a general real or rational quadratic form is needed.
These data support the next restriction comparison and do not
by themselves assert its rank.

\begin{proof}
For the determinant samples begin with the same traceless integer
matrices and generator as QC-007. If
$M(v)=v_1A_1+v_2A_2+v_3A_3$, the quadratic coefficient of
$\det(I_9+M)$ is
\[
 \sum_{a<b}(M_{aa}M_{bb}-M_{ab}M_{ba})
 =\frac{(\operatorname{tr}M)^2-\operatorname{tr}(M^2)}2
 =-\frac12\operatorname{tr}(M^2).
\]
Since the normalized second tensor is the ordinary quadratic
coefficient divided by $\binom92=36$, its symmetric matrix is
\[
 Q_{ij}=-\frac{\operatorname{tr}(A_iA_j)}{72}\pmod {101}.
 \tag{QC-009.2}
\]
The generator discards a pencil unless $Q$ is nonsingular and
$\det Q$ is a square modulo $101$.

The following finite algorithm constructs a change of basis $h$
with $h^{\mathsf T}Qh=I_3$. Apply ordinary congruence
Gram--Schmidt to the three coordinate vectors, choosing a vector
of nonzero norm at each step. If the remaining basis vectors
all have zero norms, some two have nonzero mutual pairing,
because the remaining quadratic space is nondegenerate; their
sum has nonzero norm as $2\ne0$ in the field. Orthogonal
subtraction then gives an invertible basis with diagonal Gram
matrix $\operatorname{diag}(d_1,d_2,d_3)$.

Because its determinant is a square, there are either zero or
two nonsquare $d_i$. Scale each square-norm vector by the
inverse of a square root. If $a,b$ are the two nonsquare
diagonal entries, find $x,y$ with
\[
 ax^2+by^2=1.
\]
Such a pair exists: each of the two sets
$\{ax^2:x\in\mathbb F_{101}\}$ and
$\{1-by^2:y\in\mathbb F_{101}\}$ has $51$ elements,
so they intersect. The code simply enumerates $y$ and consults
a table of all squares. Choose $r$ with $r^2=ab$, which is
possible since $a,b$ are both nonsquares. The two coefficient
vectors
\[
 (x,y),\qquad r^{-1}(-by,ax)
\]
are orthonormal for $\operatorname{diag}(a,b)$: their mutual
pairing is zero and both norms are one. This finishes the
construction of $h$. The code verifies every entry of
$h^{\mathsf T}Qh=I_3$ and verifies $\det h\ne0$.

Replace the pencil matrices by
\[
 A'_j=\sum_{i=1}^3h_{ij}A_i\pmod {101}.
 \tag{QC-009.3}
\]
For an explicit integer lift, reduce their entries modulo
$101$ and replace the last diagonal entry by minus the sum
of the first eight diagonal entries. This keeps the same
residue matrix and makes each trace zero over $\mathbb Z$.
The resulting polynomial
$\det(x_0I_9+\sum_jx_jA'_j)$ is a genuine determinant,
and its normalized second tensor is $I_3$ modulo $101$.
Exact Newton interpolation as in QC-007 computes all higher
entries. Every second-tensor entry is asserted once more
from the reconstructed coefficients, independently of
(QC-009.2). Three additional direct determinant values are
checked for every accepted point.

For the padded samples use ordinary ternary coefficient forms
$B_j$ in the expansion
\[
 F=x_0^9+\sum_{j=2}^9x_0^{9-j}B_j(x_1,x_2,x_3).
\]
Choose an integer linear form $l$, set
$B_2=36(x_1^2+x_2^2+x_3^2)$, choose arbitrary integer
$B_3,\ldots,B_8$ by the stated generator, and define
\[
 B_9=l^9+\sum_{j=2}^8(-1)^j l^{9-j}B_j.
 \tag{QC-009.4}
\]
Then $F(-l,x_1,x_2,x_3)=0$ identically over $\mathbb Z$.
An explicit quotient is obtained by $q_0=1$,
$q_1=-l$, and $q_j=B_j-lq_{j-1}$ for $2\leq j\leq8$:
\[
 F=(x_0+l)\sum_{j=0}^8q_jx_0^{8-j}.
 \tag{QC-009.5}
\]
Indeed (QC-009.4) is exactly $B_9=lq_8$. Thus this is an
integer point of the linear-factor variety, not just a
finite-field solution to sampled incidence equations. Its
normalized second tensor is $I_3$ already over $\mathbb Q$.
The code computes residues of (QC-009.4), normalizes all
tensor entries, and checks both zero remainder and the
quotient identity at three points per sample. The exact
integer lift is defined by (QC-009.4); reducing its expanded
coefficients during the modular calculation does not change
any of its recorded values.

The two files use the same ascending coordinate order
$(j,a,b,j-a-b)$ as QC-007, now for $2\leq j\leq9$.
The initial seed is $20261122$ in both cases. The first data
file records $2506$ rejected Gram matrices before obtaining
its $2420$ accepted points. All $7260$ direct determinant
checks and all $7230$ padded consistency checks passed.
The runtimes were respectively $35.062$ and $3.937$ seconds.
The files are
\texttt{orthonormal\_determinant\_2420.json} and
\texttt{orthonormal\_padded\_2410.json}, with adjacent
manifests recording source and file hashes. Their normalized
tensor-byte SHA-256 digests are respectively
\begin{quote}\ttfamily\small
f30ffaec897ba3ddca20d9a7944477027f01f49da97fe0d5df307d81f0c24d82,\\
9634a294285e7b6d333c0507bd4ffe6da9c588424daf0971d171e3e3cff103ab.
\end{quote}
The first source pencil or padded source forms, all accepted
generator states, and all consistency values are retained.
Every point is reproducible from the script and seed.
\end{proof}

\paragraph{Use of the modular slice.}
A polynomial defined by a universal bracket construction may
be simplified at $C_2=I_3$ before evaluation. At these sample
points the simplification gives exactly its original value
modulo $101$. A nonzero minor obtained this way still lifts
to a nonzero minor of the original polynomials on their
integer source points, because all normalization denominators
are units modulo $101$. This argument does not require the
chosen integer lifts of the determinant matrices to have
$C_2=I_3$ in characteristic zero.
